\section{clasificador KNN}
La clasificación por vecino más cercano representa uno de los enfoques más sencillos y, a la
vez, eficaces del aprendizaje automático. A diferencia de los algoritmos que construyen modelos
complejos durante el entrenamiento, los métodos de vecino más cercano adoptan un enfoque más
directo para la clasificación. El proceso de entrenamiento de los clasificadores por vecino más cercano
(también denominados «basados en instancias» o «basados en memoria») es notablemente sencillo:
3 consiste simplemente en almacenar todos los ejemplos de entrenamiento en una estructura de
datos adecuada para su posterior recuperación. Esto contrasta marcadamente con otros algoritmos
de clasificación que crean modelos abstractos durante el entrenamiento. Lo que hace único al método
del vecino más cercano es que pospone la mayor parte del trabajo computacional hasta el momento
de la clasificación. Cuando un nuevo ejemplo necesita ser clasificado, el algoritmo identifica los
k ejemplos de entrenamiento más similares a la nueva instancia; estos son sus «k vecinos más
cercanos». A continuación, se determina la clase del nuevo ejemplo basándose en estos vecinos. El
enfoque básico se limita a contar la clase mayoritaria entre los vecinos, pero las implementaciones
más sofisticadas utilizan la ponderación por distancia. Con la ponderación por distancia, los vecinos
físicamente más cercanos al ejemplo objetivo tienen mayor influencia en la decisión de clasificación.
En nuestro ejemplo del correo electrónico, si el único vecino «no spam» estuviera mucho más
cerca de nuestro nuevo correo electrónico que los dos vecinos «spam», un enfoque ponderado por
distancia podría seguir clasificando el nuevo correo electrónico como «no spam». La selección del
valor adecuado para k depende del contexto: no existe un número óptimo universal. Con k = 1,
la clasificación se basa únicamente en el único ejemplo más cercano, lo que podría ser sensible al
ruido. Con valores de k más altos, el algoritmo tiene en cuenta más vecinos, pero puede difuminar
los límites entre clases. Para nuestro ejemplo del correo electrónico, un conjunto de datos muy
pequeño podría funcionar mejor con k = 3, mientras que un conjunto de datos más grande podría
beneficiarse de k = 15 o superior
